\documentclass{article}
\usepackage[T1]{fontenc}
\usepackage{lmodern}
\usepackage{graphicx}
\usepackage{amsmath,amssymb}
\usepackage[a4paper,margin=2cm]{geometry}
\usepackage[absolute,overlay]{textpos}
\setlength{\TPHorizModule}{1mm}
\setlength{\TPVertModule}{1mm}
\usepackage[indonesian]{babel}
\usepackage{hyperref}
\setlength{\emergencystretch}{2em}
% unit: Halaman 20

\begin{document}

% === Halaman 20 (llm) ===

header for SPIE use

\begin{center}
\textbf{BPCS-Steganography Experimental Program site:}\
\url{http://www.datahide.org/BPCS e:QPtechHV-program-e.html}

\Large \textbf{Principle and applications of BPCS-Steganography}

\large Eiji Kawaguchi$^*$ and Richard O. Eason$^{**}$

\small $^*$ Kyushu Institute of Technology, Kitakyushu, Japan\
$^{**}$ University of Maine, Orono, Maine 04469-5708
\end{center}

\begin{center}
\textbf{ABSTRACT}
\end{center}

Steganography is a technique to hide secret information in some other data (we call it a \textit{vessel}) without leaving any apparent evidence of data alteration. All of the traditional steganographic techniques have limited information-hiding capacity. They can hide only $10\%$ (or less) of the data amounts of the vessel. This is because the principle of those techniques was either to replace a special part of the frequency components of the vessel image, or to replace all the least significant bits of a multi-valued image with the secret information.

Our new steganography uses an image as the vessel data, and we embed secret information in the bit-planes of the vessel. This technique makes use of the characteristics of the human vision system where by herself a human cannot perceive any shape information in a very complicated binary pattern. We can replace all of the ``noise-like'' regions in the bit-planes of the vessel image with secret data without deteriorating the image quality. We termed our steganography ``BPCS-Steganography,'' which stands for Bit-Plane Complexity Segmentation-steganography.

We made an experimental system to investigate this technique in depth. The merits of BPCS-Steganography found by the experiments are as follows.

\begin{enumerate}
    \item The information hiding capacity of a true color image is around $50\%$.
    \item A sharpening operation on the dummy image increases the embedding capacity quite a bit.
    \item Canonical Gray coded bit planes are more suitable for BPCS-Steganography than the standard binary bit planes.
\end{enumerate}

\end{document}
